\subsection{Volume work and finite field equations of state}
\label{sec:finite-field-equation-of-state}

An equation of state needs more than an update at fixed geometry. Energy,
volume work and state preparation must refer to the same dynamical system.
The distinction can be made exact for scalar pair repair, the golden scalar
action, and the finite Maxwell action. These statements concern bounded
self-reading patches with local state, ports, retained readback and repair
or feedback, together with their declared field encodings. Thermodynamic
preparation and metric deformation are additional inputs.

\paragraph{Fixed-volume repair does not determine pressure.}
For $Q(x)=\frac12\sum_i x_i^2$, pair averaging at $(i,j)$ preserves
$\sum_i x_i$ and decreases $Q$ by $(x_i-x_j)^2/4$. An added cumulative
loss record $B$ therefore preserves $Q+B$. This identity does not identify
the loss as heat or $Q$ as physical energy.

\begin{proposition}[Nonuniqueness of an unprovided volume response]
\label{prop:finite-eos-work-ambiguity}
Fix $V_0>0$. For every real $\alpha$, the positive extension
\[
 H_\alpha(x,V)=Q(x)(V/V_0)^{-\alpha},\qquad Q(x)>0,\quad V>0,
\]
agrees with $Q$ on every state at $V=V_0$. Under the supplied convention
$p=-\partial_VH|_x$ and $\rho=H/V$, it has $p/\rho=\alpha$.
Consequently complete fixed-volume repair histories and the function $Q$
do not select a work ratio within this extension class.
\end{proposition}
\begin{proof}
At $V_0$ the multiplicative factor is one, independently of $x$.
Differentiation gives $-\partial_VH_\alpha=\alpha H_\alpha/V$.
In particular $\alpha=1/3$ and $\alpha=1/2$ give different positive
pressures despite identical scalar histories at the reference volume.
\end{proof}
Adding a volume-independent constant to $H$ also changes $\rho$ while
leaving its work derivative unchanged. The energy reference is therefore
part of the required specification. This is neither an equilibrium
construction nor an obstruction to supplying an independently justified
volume response. Exact twelve-port traces retain both quadratic loss and
the added bookkeeping record; the latter is an external recorder, not a
derived thermal bath.

\paragraph{Stress from the golden scalar action.}
Let $s>0$ dilate the reference cube, with $V=s^3$ in fixed reference units.
The dual mass and gradient forms scale as $M_s=s^3M_0$ and
$K_s^{\rm grad}=sK_0^{\rm grad}$. Keep the mass parameter $m$ fixed in
those units. For original field coordinates $q$ and canonical momenta $p$,
define at $s=1$
\[
 K=\tfrac12p^TM_0^{-1}p,\qquad
 G=\tfrac12q^TK_0^{\rm grad}q,\qquad
 U=\tfrac12m^2q^TM_0q.
\]
The full box volume is distinct from the sum of interior dual-cell masses.

\begin{proposition}[Canonical scalar work and thermal specialization]
\label{prop:finite-eos-scalar}
At fixed original $(q,p)$,
\[
 H_s=K/s^3+sG+s^3U,\qquad
 p_{\rm sc}V=K-G/3-U,\qquad \rho_{\rm sc}=(K+G+U)/V,
\]
where the last two expressions use energies evaluated at the chosen scale.
For nonnegative $K,G,U$ and positive total energy, the instantaneous ratio
lies in $[-1,1]$. If $K=G+U$, its value lies in $[0,1/3]$; for $U=0$
it equals $1/3$.
\end{proposition}
\begin{proof}
Differentiate $H_s$ and divide by $dV/ds=3s^2$. Positivity gives the
first bounds. Substitution of $K=G+U$ gives $p_{\rm sc}V=2G/3$ and
$H=2(G+U)$, proving the remaining assertions.
\end{proof}
Holding mass-normalized canonical coordinates fixed during dilation would
instead give $pV=2G/3$. Its difference from the original-coordinate stress
is $K-G-U=\frac12d(q^Tp)/dt$ under the continuous action. Stationarity or
a controlled virial average removes this difference; transient records do
not satisfy that condition automatically.

For the authenticated $q=5$ execution, the centered velocity is
$v_n=(q_n-q_{n-1})/\tau-\tau\mathsf A q_n/2$ and $p_n=M_0v_n$.
Exact arithmetic evaluates $K,G,U$ and stress on all twenty-two retained
states of all four histories. The conserved split quantity is
$H-\tau^2\|\mathsf A q\|_{M_0}^2/8$, not $H$.
Full-operator sine and cosine Taylor polynomials, with outward rational
remainders, separately enclose the continuous-action pressure at all
twenty-two times. In detail, $\operatorname{spec}_{M_0}(\mathsf A)
\subset[1,614]$, $\|v_0\|_{M_0}\le4$ and $t\le3/2$ give the
displacement and velocity remainder bounds
\[
 \delta_q=\frac{4(3/2)^{163}614^{81}}{163!},\qquad
 \delta_v=\frac{4(3/2)^{162}614^{81}}{162!}.
\]
These follow from the real sine and cosine Taylor remainders by spectral
calculus, retaining every mode. The corresponding polynomial norms are
bounded by $7$ and $5$. Expanding the stress quadratic form therefore gives
\[
 |\delta p|\le5\delta_v+\delta_v^2/2
       +\frac{616}{3}(7\delta_q+\delta_q^2/2)<1.385\cdot10^{-33}.
\]
Outward rational rounding on the $10^{-12}$ reporting grid is separate
from this analytic tail. At $t=3/\sqrt5$ and $V=1$,
\begin{align*}
 p_{\rm continuous}&\in[0.129161530903,0.129161530906],\\
 (p/\rho)_{\rm continuous}&\in[0.467230205001,0.467230205004],\\
 p_{\rm split}-p_{\rm continuous}
   &\in[-0.017284266362,-0.017284266359].
\end{align*}
These are certified nonequilibrium stress calculations in the supplied model. The
ratio exceeding $1/3$ does not contradict its thermal specialization.

For a separately supplied canonical Bose state, let $\lambda_j>0$ be
the complete spatial generalized spectrum of the same Dirichlet action.
In fixed units $\hbar=k_B=c=1$, with zero chemical potential and zero-point
energy subtracted at each $s$, put
\[
 \omega_j^2=m^2+\lambda_j/s^2,\quad
 n_j=(e^{\omega_j/T}-1)^{-1},\quad E_j=\omega_jn_j,
 \qquad F=T\sum_j\log(1-e^{-\omega_j/T}).
\]
Differentiating this free energy at fixed $T,m$ and mode count gives
\[
 p=-\partial_VF
  =\frac1{3V}\sum_j\frac{E_j\lambda_j}{s^2\omega_j^2},\qquad
 \rho=\frac1V\sum_jE_j.
\]
The vacuum-subtracted thermal covariances obey $K=G+U$, so this pressure agrees with the
canonical mechanical stress. The massive ratio depends on temperature
and scale; the massless ratio is $1/3$ from the declared three-dimensional
scaling. All $64$ modes are retained in $36$ numerical cases with
$s\in\{1,2,8\}$, $m\in\{0,1,4\}$ and $T\in\{1/2,2,10,50\}$.
Their ratios range from approximately $0.0409561$ to $1/3$.
An independent high-precision spectrum and partition derivative check these
numerical values; this thermal calculation is not an interval certificate.
At fixed mode count the ultraviolet scale changes with $s$, so neither a
bulk thermodynamic limit nor cutoff removal follows. The Bose occupation
support is unbounded despite the finite mode count. Native preparation,
thermalization, finite observer-memory realization and physical calibration
are not supplied. Subtracting zero-point energy at every scale also
subtracts its pressure; no vacuum-pressure prediction is asserted.

\paragraph{Maxwell field pressure on the supplied cone.}
For the declared three-dimensional Whitney action, let $a$ be the
edge-integrated potential and $\pi$ its conjugate momentum. Under uniform
metric dilation, the one-form mass is $M_s=sM_1$ and the two-form mass is
$N_s=s^{-1}N_1$. With curvature $b=Ca$,
\[
 H_s=\tfrac12\pi^TM_s^{-1}\pi+\tfrac12b^TN_sb=H_1/s.
\]
At fixed $(a,\pi)$, $V_s=s^3V_1$ therefore gives
$p_{\rm mean}=H_s/(3V_s)$ for positive field energy. Independently the
integrated Maxwell momentum-flux tensor
\[
 \Pi=\frac1V\int\left[
 \tfrac12(|E|^2+|B|^2)I-EE^T-BB^T\right]d^3x
\]
has $\operatorname{tr}\Pi/3=H/(3V)$. Directional stresses can include
tension; the trace identity does not require an isotropic individual field.
Zero field energy leaves the ratio undefined. This is field energy only,
excluding matter, sources, clocks, walls and observer records.

On the complete cone, removing its twelve gradient gauge directions
leaves thirty positive oscillator modes. A separately supplied classical
Gibbs ensemble on a fixed canonical gauge quotient, with a geometry-independent
Liouville cell normalization, has
\[
 \log Z=30\log T-\sum_{j=1}^{30}\log\omega_j+\text{constant},
 \qquad \langle H\rangle=30T,\qquad pV=10T.
\]
The formulas follow by the finite Gaussian integral and
$\omega_j(s)=\omega_j(1)/s$. Independent gauge quotients and spatial
quadrature check nine scale/temperature cases. A separate source-free
numerical episode retains eighty nonzero displacement-current updates and
eighty-one full canonical states. It checks energy balance, Amp\`ere,
Faraday and Gauss with zero conduction current. The pressure extension also
consumes the decoded fields of the self-reading parent instrument; the new
source-free trajectory is not itself another observer readout instrument.
The Gibbs preparation is separate from that trajectory, and supplies no
thermalization, quantum blackbody law or continuum limit.

The work nonuniqueness and finite scalar/Maxwell scaling identities are
kernel-checked in \path{Lean/Thermodynamics/EquationOfState.lean}.
The state, stress and independent verification artifacts are in
\path{code/equation_of_state/}. Gaussian ensemble evaluation and Maxwell
trajectory checks use numerical tolerances; only the scalar continuous
comparison has the stated outward-rational pressure enclosures.
These are consequences of supplied action and state laws, not a derivation
of a thermodynamic equation of state from scalar pair averaging.
